% SEGMENT OLP-0553-S01
\documentclass[../../../include/open-logic-section]{subfiles}

\begin{document}

\olfileid{sth}{ordinals}{basic}
\olsection{வரிசையெண்களின் அடிப்படைப் பண்புகள்}

முதல் சில வரிசையெண்கள் இயல் எண்களே என்று கண்டோம். இயல் எண்கள்மீது
செல்லுபடியாகும் தொகுத்தறிதல் கொள்கையை நீட்டிப்பதே வரிசையெண்களின் கோட்பாட்டை
உருவாக்குவதற்கான முதன்மைக் காரணம். தொடக்கநிலை முடிவுகளின் ஒரு தொடரின்
வழியாக இதை நோக்கி முன்னேறுவோம்.

\begin{lem}\ollabel{ordmemberord}
ஒரு வரிசையெண்ணின் ஒவ்வோர் !!{element} உம் ஒரு வரிசையெண்.
\end{lem}

\begin{proof}
$\alpha$ ஒரு வரிசையெண் என்றும் $b \in \alpha$ என்றும் கொள்வோம்.
$\alpha$ கடப்புப் பண்புடையதால் $b \subseteq \alpha$. ஆகவே $\in$ ஆனது
$\alpha$ ஐ நன்கு வரிசைப்படுத்துவதால், $\in$ ஆனது $b$ ஐயும் நன்கு
வரிசைப்படுத்துகிறது.

$b$ கடப்புப் பண்புடையது என்று காண, $x \in c \in b$ எனக் கொள்வோம்.
$b \subseteq \alpha$ என்பதால் $c \in \alpha$. மீண்டும் $\alpha$
கடப்புப் பண்புடையதால் $c \subseteq \alpha$; எனவே $x \in \alpha$.
ஆகவே $x, c, b \in \alpha$. ஆனால் $\in$ ஆனது $\alpha$ ஐ நன்கு
வரிசைப்படுத்துவதால், \olref[sth][ordinals][wo]{wo:strictorder} இன்படி
$\in$ ஆனது $\alpha$ மீதான கடப்புப் பண்புடைய தொடர்பு. எனவே $x \in c
\in b$ என்பதிலிருந்து $x \in b$. பொதுமைப்படுத்தினால், $c \subseteq b$.
\end{proof}

\begin{cor}\ollabel{ordissetofsmallerord}
எந்த வரிசையெண்~$\alpha$ க்கும்,
$\alpha = \Setabs{\beta \in \alpha}{\beta \text{ ஒரு வரிசையெண்}}$.
\end{cor}

\begin{proof}
\olref{ordmemberord} இலிருந்து உடனடியாகப் பின்வருகிறது.
\end{proof}

அடுத்த இரண்டு முதன்மை முடிவுகளான \olref{ordinductionschema},
\olref{ordtrichotomy} ஆகியவற்றின் தோராயமான கருத்து, வரிசையெண்களையே
உறுப்புறவு நன்கு வரிசைப்படுத்துகிறது என்பதாகும்:

\begin{thm}[முடிவிலிக்கடந்த தொகுத்தறிதல்]\ollabel{ordinductionschema}
எந்த வாய்பாடு $\phi(x)$ க்கும்:	
\[
	\text{எனில் }\exists \alpha \phi(\alpha)\text{, அப்போது }\exists \alpha(\phi(\alpha)
	\land  (\forall \beta \in \alpha) \lnot \phi(\beta))
\]
இங்கே காட்டப்பட்ட அளவீடுகள் வரிசையெண்களுக்கு மறைமுகமாகக்
கட்டுப்படுத்தப்பட்டுள்ளன.
\end{thm}

\begin{proof}
%\olref{ordinductionschema1}. 
ஏதோ ஒரு வரிசையெண் $\alpha$ க்கு $\phi(\alpha)$ எனக் கொள்வோம்.
$ (\forall \beta \in \alpha) \lnot \phi(\beta)$ எனில், நிறுவல்
முடிந்தது. இல்லையெனில், $\alpha$ ஒரு வரிசையெண் என்பதால், அதில்
$\phi$-ஐ நிறைவுசெய்யும் உறுப்புகளில் ஓர் $\in$-மீச்சிறு !!{element}
உண்டு; \olref{ordmemberord} இன்படி இதுவும் ஒரு வரிசையெண்.
% SOURCE CORRECTION TA-STH-011: the least element satisfies phi; it is not “which is phi.”
%	\olref{ordinductionschema2}. Suppose there is some ordinal
%	$\gamma$ such that $\lnot\phi(\gamma)$. Then by
%	\olref{ordinductionschema1}, there is an $\in$-minimal ordinal
%	$\alpha$ for which $\lnot\phi(\alpha)$. So $(\forall \beta <
%	\alpha) \phi(\beta)$, rendering the antecedent of the conditional
%	false.
\end{proof}
\noindent
\olref{ordinductionschema} ஐப் பின்வரும் திட்டமாகவும் சமமாக
வெளிப்படுத்தலாம் என்பதைக் கவனியுங்கள்:
\[
\text{எனில் }\forall \alpha((\forall \beta \in \alpha)\phi(\beta) \lif 
\phi(\alpha))\text{, அப்போது }\forall \alpha\phi(\alpha)
\]
இதற்கு \olref{ordinductionschema} இல் $\lnot\phi(\alpha)$ ஐ எடுத்துக்கொண்டு,
பின்னர் தொடக்கநிலைத் தருக்கச் செயல்முறைகளைச் செய்தால் போதும்.

\begin{thm}[முக்கூற்றுப் பண்பு]\ollabel{ordtrichotomy}
எந்த வரிசையெண்கள் $\alpha$, $\beta$ க்கும்
$\alpha \in \beta \lor \alpha = \beta \lor \beta \in \alpha$.
\end{thm}

\begin{proof}
இது இரட்டைத் தொகுத்தறிதல் நிறுவல்; அதாவது
\olref{ordinductionschema} ஐ இருமுறை பயன்படுத்துகிறோம். $x \in y \lor
x = y \lor y \in x$ எனில், அப்போது, அப்போது மட்டுமே, $x$ ஆனது $y$
உடன் \emph{ஒப்பிடத்தக்கது} என்று கூறுவோம்.

தொகுத்தறிதலுக்காக, $\alpha$ விலுள்ள ஒவ்வொரு வரிசையெண்ணும்
\emph{ஒவ்வொரு} வரிசையெண்ணுடனும் ஒப்பிடத்தக்கது எனக் கொள்வோம். மேலும் ஒரு
தொகுத்தறிதலுக்காக, $\beta$ விலுள்ள ஒவ்வொரு வரிசையெண்ணுடனும் $\alpha$
ஒப்பிடத்தக்கது எனக் கொள்வோம். $\alpha$ ஆனது $\beta$ உடன் ஒப்பிடத்தக்கது
என்று காட்டுவோம். $\beta$ மீதான தொகுத்தறிதலால், $\alpha$ ஒவ்வொரு
வரிசையெண்ணுடனும் ஒப்பிடத்தக்கது என்று பின்வரும்; அதன்பின் $\alpha$ மீதான
தொகுத்தறிதலால், தேவைப்பட்டவாறு, \emph{ஒவ்வொரு} வரிசையெண்ணும்
\emph{ஒவ்வொரு} வரிசையெண்ணுடனும் ஒப்பிடத்தக்கது என்று பின்வரும்.
$\alpha \notin \beta$, $\beta \notin \alpha$ என்று கொண்டு,
$\alpha = \beta$ என்று காட்டினால் போதும்.

$\alpha \subseteq \beta$ என்று காட்ட, $\gamma \in \alpha$ ஐ
நிலைப்படுத்துங்கள்; \olref{ordmemberord} இன்படி இது ஒரு வரிசையெண்.
ஆகவே முதல் தொகுத்தறிதல் கருதுகோளின்படி, $\gamma$ ஆனது $\beta$ உடன்
ஒப்பிடத்தக்கது. ஆனால் $\gamma = \beta$ அல்லது $\beta \in \gamma$
என்ற இரண்டில் எது உண்மையானாலும், $\alpha$ கடப்புப் பண்புடையது என்பதைத்
தேவைப்பட்டால் பயன்படுத்தி, $\beta \in \alpha$ என்று பின்வரும்; இது நமது
கருதுகோளுக்கு முரணானது. எனவே $\gamma \in \beta$. பொதுமைப்படுத்தினால்,
$\alpha \subseteq \beta$.

இரண்டாவது தொகுத்தறிதல் கருதுகோளைப் பயன்படுத்தும் முற்றிலும் இதேபோன்ற
பகுத்தறிதல், $\beta \subseteq \alpha$ என்று காட்டுகிறது. எனவே
$\alpha = \beta$.
\end{proof}\noindent
இதனால், $\in$ ஒரு வரிசைப்படுத்தும் தொடர்பாகச் செயல்படுவதால்,
$\alpha \in \beta$ என்பதற்குப் பதிலாகச் சில நேரங்களில் $\alpha <\beta$
என்று எழுதுவோம். பழக்கம் என்பதையும், $\alpha \in \beta \lor \alpha =
\beta$ என்று எழுதுவதைவிட $\alpha \leq \beta$ என்று எழுதுவது எளிது
என்பதையும் தவிர, இதற்கு ஆழமான காரணங்கள் இல்லை.\footnote{$\alpha
\mathrel{\underline{\in}} \beta$ என்று எழுதலாம்; ஆனால் அது முற்றிலும்
தரநிலையற்றதாக இருக்கும்.}

நமது கடைசி முடிவுகளிலிருந்து விரைவாகப் பின்வரும் இரண்டு விளைவுகள் இதோ;
முதலாவது நமது புதிய குறியீட்டைப் பயன்படுத்துகிறது:

\begin{cor}\ollabel{ordordered}
$\exists \alpha\phi(\alpha)$ எனில், $\exists \alpha(\phi(\alpha)
\land \forall \beta(\phi(\beta) \lif \alpha \leq \beta))$.
மேலும், எந்த வரிசையெண்கள் $\alpha, \beta, \gamma$ க்கும்,
$\alpha \notin \alpha$ மற்றும் $\alpha \in \beta \in \gamma \lif
\alpha \in \gamma$ ஆகிய இரண்டும் உண்மை.
\end{cor}

\begin{proof}
\olref[wo]{wo:strictorder} போலவே.
\end{proof}

\begin{prob}
\olref[sth][ordinals][basic]{ordtrichotomy} இன் நிறுவலில் உள்ள
``முற்றிலும் இதேபோன்ற பகுத்தறிதலை'' நிறைவுசெய்க.
\end{prob}

\begin{cor}\ollabel{corordtransitiveord}
$A$ ஒரு வரிசையெண் எனில், எனில் மட்டுமே, $A$ வரிசையெண்களைக் கொண்ட
கடப்புப் பண்புடைய கணம்.
\end{cor}

\begin{proof}
\emph{இடமிருந்து வலம்.} \olref{ordmemberord} இன்படி.
\emph{வலமிருந்து இடம்.} $A$ வரிசையெண்களைக் கொண்ட கடப்புப் பண்புடைய
கணம் எனில், \olref{ordinductionschema}, \olref{ordtrichotomy}
ஆகியவற்றின்படி $\in$ ஆனது $A$ ஐ நன்கு வரிசைப்படுத்துகிறது.
\end{proof}

\olref{ordinductionschema}, \olref{ordtrichotomy} ஆகியவை $\in$ ஆனது
வரிசையெண்களை நன்கு வரிசைப்படுத்துகிறது என்று கூறுவதாக நாம் சுருக்கினோம்.
ஆனால் பின்வரும் முடிவின் காரணமாக, இவ்வகைக் கூற்றைப் பற்றி நாம்
\emph{மிகவும் எச்சரிக்கையாக} இருக்க வேண்டும்:

\begin{thm}[புராலி--ஃபோர்ட்டி முரணுரை]\ollabel{buraliforti}
எல்லா வரிசையெண்களையும் கொண்ட கணம் எதுவும் இல்லை.
\end{thm}

\begin{proof}
முரண் மூலம் நிறுவும் பொருட்டு, $O$ எல்லா வரிசையெண்களையும் கொண்ட கணம்
எனக் கொள்வோம். $\alpha \in \beta \in O$ எனில்,
\olref{ordmemberord} இன்படி $\alpha$ ஒரு வரிசையெண்; ஆகவே
$\alpha \in O$. எனவே $O$ கடப்புப் பண்புடையது; இதனால்
\olref{corordtransitiveord} இன்படி $O$ ஒரு வரிசையெண். ஆகவே $O \in O$;
இது \olref{ordordered} உடன் முரண்படுகிறது.
\end{proof}
\noindent
இந்த முடிவு \citeauthor{Burali-Forti1897} என்பவரின் பெயரால்
அழைக்கப்படுகிறது. ஆனால் 1899 இல் டெடெகிண்டிற்கு எழுதிய கடிதத்தில்,
எல்லா வரிசையெண்களையும் கொண்ட ஒரு கணம் உள்ளது என்ற அனுமானத்திலுள்ள
\emph{முரண்பாட்டை} முதன்முதலில் தெளிவாகக் கண்டவர் காண்டரே. ஃபான்
ஹெயெனூர்ட் விளக்குவது போல:
\begin{quote}
  வரிசையெண்களின் இயல்பான வரிசை ஒரு பகுதி வரிசை மட்டுமே என்ற முடிவை,
  \emph{முரண்வழி நிறுவலால்}, நிறுவுவதாகவே புராலி--ஃபோர்ட்டி அந்த
  முரண்பாட்டைக் கருதினார்.
  \citep[p.~105]{Heijenoort1967}
\end{quote}
புராலி--ஃபோர்ட்டியின் தவறை ஒதுக்கிவிட்டு, மேற்கண்டவற்றைப் பின்வருமாறு
சுருக்கலாம். வரிசையெண்கள், உறுப்புறவால் தனித்தனியாக நன்கு
வரிசைப்படுத்தப்பட்ட கணங்கள்; ஒரு கணத்தைத் தொகுப்பாக அமைக்காமலேயே,
உறுப்புறவால் ஒட்டுமொத்தமாகவும் நன்கு வரிசைப்படுத்தப்பட்டவை.

இதை நிறைவுசெய்ய, \olref{ordinductionschema},
\olref{ordtrichotomy} ஆகியவற்றிலிருந்து பின்வரும் வரிசையெண்களின் இன்னும்
சில அடிப்படைப் பண்புகளைப் பார்ப்போம்.

\begin{prop}
வரிசையெண்களைக் கொண்ட கண்டிப்பாக இறங்கும் எந்தத் தொடரும் முடிவுறு தொடராகும்.
\end{prop}

\begin{proof}
வரிசையெண்களைக் கொண்ட முடிவுறாத கண்டிப்பாக இறங்கும் எந்தத் தொடருக்கும்
$\alpha_0 > \alpha_1 > \alpha_2 > \ldots$ ஒரு $<$-குறும உறுப்பு
இல்லை; இது \olref{ordinductionschema} உடன் முரண்படுகிறது.
\end{proof}

\begin{prop}\ollabel{ordinalsaresubsets}
எந்த வரிசையெண்கள் $\alpha, \beta$ க்கும்
$\alpha \subseteq \beta \lor \beta \subseteq \alpha$.
\end{prop}

\begin{proof}
$\alpha \in \beta$ எனில், $\beta$ கடப்புப் பண்புடையதால்
$\alpha \subseteq \beta$. இதேபோல, $\beta \in \alpha$ எனில்
$\beta \subseteq \alpha$. மேலும் $\alpha = \beta$ எனில்,
$\alpha \subseteq \beta$, $\beta \subseteq \alpha$ ஆகிய இரண்டும்
உண்மை. ஆகவே \olref{ordtrichotomy} இன்படி நிறுவல் முடிந்தது.
\end{proof}

\begin{prop}\ollabel{ordisoidentity}
எந்த வரிசையெண்கள் $\alpha, \beta$ க்கும், $\alpha = \beta$ எனில்,
எனில் மட்டுமே, $\ordeq{\alpha}{\beta}$.
\end{prop}

\begin{proof}
வரிசையெண்கள் நன்கு வரிசைப்படுத்தல்கள்; எனவே முக்கூற்றுப் பண்பு
(\olref[basic]{ordtrichotomy}), \olref[iso]{wellordnotinitial}
ஆகியவற்றிலிருந்து இது உடனே பின்வருகிறது.
\end{proof}

\begin{prob}\ollabel{probunionordinalsordinal}
$X$ இன் ஒவ்வோர் உறுப்பும் ஒரு வரிசையெண் எனில், $\bigcup X$ ஒரு
வரிசையெண் என்று நிறுவுக.
\end{prob}

\end{document}
